\chapter{Metodologia di lavoro e selezione dei criteri WCAG}
\label{chap:metodologia-selezione}

Questo capitolo tratta la definizione del metodo di lavoro adottato e la selezione del sottoinsieme di criteri WCAG~2.1 livello A e AA su cui è stata concentrata la verifica automatica. Le fasi successive del lavoro (progettazione e sviluppo del prototipo (capitolo~\ref{chap:progettazione-prototipo}) e validazione sperimentale (capitolo~\ref{chap:validazione-benchmark})) si basano direttamente sui criteri individuati in questo capitolo.

\section{Metodo di lavoro}
\label{sec:metodo-lavoro}

\todoplaceholder{pianificazione, interazioni con il tutor, strumenti di tracciamento.}

% TODO: pianificazione delle attività, modalità di interazione con il tutor aziendale,
% cadenza delle revisioni di progresso, strumenti usati per tracciare requisiti e progressi.

\section{Selezione dei criteri WCAG oggetto di verifica}

\label{sec:selezione-criteri}

La selezione del sottoinsieme di criteri di successo WCAG~2.1 di livello A e AA da sottoporre a verifica automatica è stata condotta a partire dall'insieme dei criteri applicabili alle pagine web, attraverso un percorso articolato in cinque fasi:

\begin{enumerate}

\item individuazione dell'insieme completo dei 50 criteri applicabili alle pagine web;

\item classificazione dei criteri in base al principio POUR e al livello di conformità;

\item verifica della presenza dei criteri nei tre report di riferimento, assunti come indicatori della frequenza con cui le diverse problematiche di accessibilità emergono sul web reale: il WebAIM Million\footcite{webaim:million2026}, il monitoraggio semplificato\footcite{agid:errori-semplificato} e il monitoraggio approfondito\footcite{agid:errori-approfondito} condotti da AGID;

\item analisi dei soli criteri riscontrati in almeno un report, considerando la verificabilità statica, la classificazione, l'integrabilità con LLM, le tecniche sufficienti WCAG di riferimento\footcite[Sono le \emph{Sufficient Techniques} ufficiali del W3C nelle quali ogni criterio ha un elenco di metodi concreti già collaudati per soddisfarlo, ciascuno con un codice (lettera = tecnologia: H=HTML, ARIA=ARIA, C=CSS, G=General) + numero, e una pagina propria che spiega esattamente come implementarlo.][]{w3c:techniques} e le condizioni d'uso coperte;

\item selezione finale dei criteri \textbf{core}, \textbf{opzionali} ed \textbf{esclusi}, rapportata alla fattibilità della verifica automatica e al tempo disponibile per la realizzazione del progetto.
\end{enumerate}

\subsection{I 50 criteri di riferimento, per principio e livello di conformità}

Come richiamato nella sezione~\ref{sec:livello-aa}, i criteri di successo di livello A e AA delle WCAG~2.1 applicabili alle pagine web sono complessivamente 50. La tabella~\ref{tab:distribuzione-pour} ne riporta la distribuzione per principio POUR (v.~sezione~\ref{sec:quattro-principi}) e per livello di conformità.

\begin{table}[htbp]
\centering
\small
\caption{Distribuzione dei 50 criteri di successo di livello A e AA per principio POUR e livello di conformità.}
\label{tab:distribuzione-pour}
\begin{tabular}{|l|r|r|r|}
\hline
\textbf{Principio} & \textbf{Livello A} & \textbf{Livello AA} & \textbf{Totale} \\
\hline
Percepibile & 9 & 11 & 20 \\
\hline
Utilizzabile & 14 & 3 & 17 \\
\hline
Comprensibile & 5 & 5 & 10 \\
\hline
Robusto & 2 & 1 & 3 \\
\hline
\textbf{Totale} & \textbf{30} & \textbf{20} & \textbf{50} \\
\hline
\end{tabular}
\end{table}
\FloatBarrier

\subsection{Copertura dei criteri nei tre report di riferimento}

A partire da questo insieme di 50 criteri, il passo successivo consiste nell'individuare quali di essi compaiono effettivamente, come categoria di errore rilevata, in almeno uno dei tre report assunti come indicatori di frequenza sul web reale: 

\begin{itemize}
\item il \textbf{WebAIM Million}, già introdotto nella sezione~\ref{sec:problema-pratico}, che analizza mediante il motore automatico WAVE le \emph{home page} di un milione di siti \emph{web} e riporta, per l'edizione 2026, la percentuale di \emph{home page} in cui compaiono le categorie di errore più comuni\footcite{webaim:million2026}; 
\item il \textbf{monitoraggio semplificato} di \gls{agid-descrizione}, condotto automaticamente da \gls{mauve-descrizione} sui siti web della Pubblica Amministrazione\footcite{agid:errori-semplificato}; 
\item il \textbf{monitoraggio approfondito} di AGID, condotto invece da un gruppo di esperti di accessibilità\footcite{agid:errori-approfondito}. 
\end{itemize}

Per ciascun criterio è stato quindi verificato in quanti di questi tre report compaia: 3/3 se presente in tutti, 2/3 se presente solo in due o 1/3 se presente solo in un solo report. La tabella~\ref{tab:frequenza-errori} riporta il risultato di questo confronto per i 13 criteri risultati presenti in almeno un report, ordinati per grado di copertura decrescente.

\begin{table}[h]
\centering
\small
\caption{Criteri di successo riscontrati nei tre report di riferimento}
\label{tab:frequenza-errori}
\begin{tabular}{|p{1.1cm}|p{3cm}|p{1.8cm}|p{1.8cm}|p{1.8cm}|p{0.9cm}|}
\hline
\textbf{Crit.} & \textbf{Nome} & \textbf{WebAIM Million} & \textbf{AGID sempl.} & \textbf{AGID appr.} & \textbf{Cop.} \\
\hline
\multicolumn{6}{|c|}{\textbf{Presenti in tutti e tre i report (3/3)}} \\
\hline
1.1.1 & Non-text Content & 53,1\% & 9,7\% & 6,5\% & 3/3 \\
\hline
1.4.3 & Contrast (Minimum) & 83,9\% & 15,1\% & 11,4\% & 3/3 \\
\hline
2.4.4 & Link Purpose (In Context) & 46,3\% & 2,3\% & 11,1\% & 3/3 \\
\hline
4.1.2 & Name, Role, Value & 30,6\% & 1,1\% & 17,8\% & 3/3 \\
\hline
\multicolumn{6}{|c|}{\textbf{Presenti in due report (2/3)}} \\
\hline
1.3.1 & Info and Relationships & --- & 1,6\% & 23,2\% & 2/3 \\
\hline
1.4.1 & Use of Color & --- & 17,4\% & 3,5\% & 2/3 \\
\hline
2.4.7 & Focus Visible & --- & 42,4\% & 2,6\% & 2/3 \\
\hline
3.3.2 & Labels or Instructions & 51\% & --- & 6,3\% & 2/3 \\
\hline
4.1.1 & Parsing & --- & 3,5\% & 3,3\% & 2/3 \\
\hline
\multicolumn{6}{|c|}{\textbf{Presenti in un solo report (1/3)}} \\
\hline
1.4.10 & Reflow & --- & 1,5\% & --- & 1/3 \\
\hline
1.4.11 & Non-text Contrast & --- & 3,5\% & --- & 1/3 \\
\hline
2.4.1 & Bypass Blocks & --- & --- & 5,1\% & 1/3 \\
\hline
3.1.1 & Language of Page & 13,5\% & --- & --- & 1/3 \\
\hline
\end{tabular}
\end{table}
\FloatBarrier

I tre report non sono equivalenti fra loro, e questo spiega perché diversi criteri compaiano solo in una parte di essi. Infatti, il WebAIM Million e il monitoraggio semplificato AGID, essendo basati esclusivamente su motori automatici, possono riportare soltanto le problematiche rilevabili senza interazione con la pagina. Conseguentemente, un criterio come 2.4.1 (\emph{Bypass Blocks}), la cui verifica richiede di valutare se i landmark presenti coprono davvero tutte le sezioni ripetute, compare solo grazie al monitoraggio approfondito, condotto da valutatori esperti, che copre anche questi casi pur con un campione di siti più ristretto rispetto al milione di \emph{home page} del WebAIM Million. Le tre fonti sono quindi complementari: le prime due forniscono un'indicazione quantitativa su scala molto ampia mediante rilevazione automatica, la terza permette di stimare la frequenza anche dei criteri che sfuggono a una rilevazione di quel tipo.

Il confronto diretto fra le percentuali del monitoraggio semplificato e di quello approfondito evidenzia inoltre un fenomeno rilevante ai fini della selezione: per alcuni criteri la quota di errore rilevata dal controllo manuale è nettamente superiore a quella rilevata dal solo controllo automatico; per altri vale l'opposto. I due casi hanno un peso molto diverso ai fini della selezione, e vengono perciò discussi separatamente.
 
Il primo caso, manuale nettamente superiore ad automatico, è quello di 1.3.1 (\emph{Info and Relationships}), 4.1.2 (\emph{Name, Role, Value}) e 2.4.4 (\emph{Link Purpose}). Il più marcato è 1.3.1, che passa dall'1,6\% degli errori rilevati automaticamente al 23,2\% di quelli confermati dagli esperti, diventando la quota più alta fra tutti i criteri del monitoraggio approfondito, superiore anche a quella di 4.1.2 (17,8\%) e 2.4.4 (11,1\%). Questi tre criteri condividono una caratteristica, ossia la loro violazione richiede un giudizio semantico che risponda a domande come "la struttura riflette davvero il contenuto?", "il nome esposto corrisponde al comportamento reale?", "il testo del collegamento è comprensibile fuori dal proprio contesto?" a cui un motore di scanzione puramente sintattico non può rispondere. Questo primo caso costituisce quindi un argomento solido, distinto dalla semplice frequenza, a favore dell'integrazione di un LLM: i criteri per cui il controllo manuale rileva sistematicamente più errori di quello automatico sono proprio quelli in cui un sistema capace di un giudizio semantico promette il maggiore valore aggiunto rispetto agli strumenti puramente sintattici oggi disponibili.
 
Il secondo caso, automatico nettamente superiore a manuale, riguarda invece 2.4.7 (\emph{Focus Visible}) e 1.4.1 (\emph{Use of Color}), la cui quota crolla rispettivamente dal 42,4\% al 2,6\% e dal 17,4\% al 3,5\%. Un divario così ampio, nella direzione opposta al primo caso, ammette più di una spiegazione: può indicare una sovra-segnalazione da parte del controllo automatico, ma può derivare anche da altri fattori non distinguibili dai soli dati dei due report, come una diversa composizione del campione di siti fra monitoraggio semplificato e approfondito. Non disponendo di dati sui singoli casi rilevati questa spiegazione resta un'ipotesi che la validazione sperimentale del capitolo~\ref{chap:validazione-benchmark} permetterà di verificare direttamente sul benchmark costituito. A differenza del primo caso, questo secondo fenomeno non è stato quindi considerato un argomento a favore dell'inclusione fra i criteri \emph{core}.



\subsection{Analisi dei criteri riscontrati nei report}

Per ciascuno dei 13 criteri riscontrati in almeno un report è stata condotta un'analisi in termini di verificabilità a partire dal solo markup HTML statico (Sì, No, Parziale); classificazione del controllo (automatico, semi-automatico, manuale), con relativa motivazione; possibilità di integrare un LLM nel controllo (Sì, No, Parziale), con relativa spiegazione; le tecniche sufficienti WCAG di riferimento\footcite[Sono le \emph{Sufficient Techniques} ufficiali del W3C nelle quali ogni criterio ha un elenco di metodi concreti già collaudati per soddisfarlo, ciascuno con un codice (lettera = tecnologia: H=HTML, ARIA=ARIA, C=CSS, G=General) + numero, e una pagina propria che spiega esattamente come implementarlo.][]{w3c:techniques}; e le condizioni d'uso (permanenti, temporanee o situazionali) che il criterio contribuisce a coprire. 

L'analisi si basa sulla documentazione ufficiale del W3C, in particolare le pagine \emph{Understanding} e le tecniche sufficienti associate a ciascun criterio\footcite{w3c:understanding-intro,w3c:quickref}. I criteri sono presentati nello stesso ordine di copertura della tabella~\ref{tab:frequenza-errori}, e una tabella riassuntiva al termine della sezione (tabella~\ref{tab:sintesi-analisi}) ne riepiloga i risultati.

\paragraph{Criteri presenti in tutti e tre i report (3/3)}

\paragraph{1.1.1 Non-text Content (Livello A, Percepibile)}
\begin{description}
\item[Verificabile via HTML statico?] Parziale.
\item[Classificazione] Semi-automatica. La presenza o assenza dell'attributo \texttt{alt} è rilevabile dal solo markup; la sua adeguatezza (corrispondenza col contenuto reale, specificità del messaggio funzionale, distinzione fra immagine decorativa e informativa) richiede invece di valutare l'immagine e il suo contesto d'uso.
\item[Integrabile con LLM?] Sì, secondo due modalità complementari:
    \begin{enumerate}
        \item Confronto visivo-testuale, in cui un LLM multimodale confronta il contenuto reale dell'immagine con l'\texttt{alt} dichiarato, individuando testi alternativi mancanti, generici (ad esempio il nome del file, o \enquote{immagine}) o incoerenti con quanto raffigurato. 
        \item Valutazione contestuale guidata da prompt: il modello riceve l'immagine, l'\texttt{alt} dichiarato e il testo circostante nella pagina, e deduce autonomamente quale sia il messaggio funzionale che l'immagine dovrebbe veicolare in quel punto della pagina, orientato da indicazioni specifiche per tipologia di sito (ad esempio e-commerce, hôtellerie, editoria) forniti tramite prompt. Sulla base di questa inferenza il modello produce un giudizio, oppure, quando l'inferenza è incerta, domande guida per chi effettua la verifica manuale, ricalcando la logica del Decision Tree del W3C\footcite{w3c:decisional-tree}.
    \end{enumerate}
\item[Tecniche di riferimento] \emph{Using \texttt{alt} attributes on \texttt{img} elements} (H37); \emph{using \texttt{aria-label} to provide labels for objects} (ARIA6/ARIA10); \emph{using null \texttt{alt} text for decorative images} (H67); \emph{using CSS to include decorative images} (C9).
\item[Condizioni d'uso coperte] Permanenti: ciechi e ipovedenti che usano uno \emph{screen reader}. Temporanee: infortunio o bendaggio che impedisce di vedere lo schermo. Situazionali: immagini non caricate per banda scarsa, schermo con riflessi o scarsa luminosità, uso a \enquote{mani libere/occhi altrove} (ad esempio alla guida, con audio).
\end{description}

\paragraph{1.4.3 Contrast (Minimum) (Livello AA, Percepibile)}
\begin{description}
\item[Verificabile via HTML statico?] Parziale.
\item[Classificazione] Semi-automatica. Se lo sfondo è un colore CSS solido dichiarato, il contrasto è calcolabile interamente da markup e CSS statici, applicando la formula di \gls{luminanza-relativa} (soglie 4{,}5:1 per il testo normale, 3:1 per il testo grande). Se lo sfondo è un'immagine, un gradiente, o dipende da stili calcolati a runtime, serve invece il rendering effettivo per campionare il colore realmente visualizzato dietro al testo; restano inoltre da valutare le eccezioni (testo decorativo, loghi, testo incidentale in una fotografia).
\item[Integrabile con LLM?] Sì, soprattutto per le eccezioni: un LLM multimodale può giudicare se un testo è decorativo, un logo o testo incidentale dentro una foto, casi che il solo calcolo matematico non distingue, e può inoltre stimare il contrasto su sfondi complessi quando il calcolo esatto non è applicabile.
\item[Tecniche di riferimento] \emph{Ensuring a contrast ratio of at least 4.5:1} (G18); \emph{ensuring a contrast ratio of at least 3:1 for large text} (G145); non specificare colori di sfondo/testo che impediscano all'utente di cambiarli (G148); fornire un controllo per passare a una presentazione ad alto contrasto (G174).
\item[Condizioni d'uso coperte] Permanenti: ipovedenti e persone con bassa sensibilità al contrasto per patologie oculari (degenerazione maculare, cataratta). Temporanee: recupero post-operatorio oculare, fotosensibilità da emicrania. Situazionali: luce solare diretta sullo schermo, schermi di scarsa qualità o luminosità.
\end{description}

\paragraph{2.4.4 Link Purpose (In Context) (Livello A, Utilizzabile)}
\begin{description}
\item[Verificabile via HTML statico?] Parziale.
\item[Classificazione] Semi-automatica. Rilevare link vuoti o privi di testo accessibile è un controllo statico; giudicare se il testo è sufficientemente descrittivo, da solo o insieme al contesto (frase, elemento di lista, cella di tabella con relativa intestazione), richiede invece comprensione semantica.
\item[Integrabile con LLM?] Sì: è un compito puramente testuale, che non richiede capacità di visione. Un LLM legge il testo del link e il contesto in cui è inserito, e giudica se lo scopo del link è chiaro, individuando pattern come \enquote{clicca qui} o \enquote{leggi di più} ripetuti senza un contesto che li distingua.
\item[Tecniche di riferimento] \emph{Creating link text that describes the purpose of a link} (G91); \emph{providing link text that describes the purpose of a link} (H30); \emph{identifying the purpose of a link using link text combined with the enclosing sentence} (G53); uso di \texttt{aria-label}/\texttt{aria-labelledby} per link con testo ambiguo (ARIA7/ARIA8).
\item[Condizioni d'uso coperte] Permanenti: utenti che usano \emph{screen reader}, che spesso navigano per elenco di link isolati dal contesto visivo. Situazionali: scansione rapida della pagina, navigazione vocale.
\end{description}

\paragraph{4.1.2 Name, Role, Value (Livello A, Robusto)}
\begin{description}
\item[Verificabile via HTML statico?] Parziale.
\item[Classificazione] Semi-automatica. La presenza di un nome accessibile (\texttt{aria-label}, \texttt{aria-labelledby}, testo interno), la validità sintattica degli attributi ARIA e l'assegnazione dei ruoli sono rilevabili dal markup; la corretta corrispondenza semantica tra ruolo assegnato e comportamento del widget e la corretta trasmissione delle notifiche di cambiamento di stato richiedono, invece, l'esecuzione o l'interazione con una tecnologia assistiva.
\item[Integrabile con LLM?] Sì: un LLM può confrontare il comportamento e l'aspetto visivo di un componente (ad esempio un elemento \texttt{div} con gestore di click che si comporta come un pulsante) con il ruolo ARIA effettivamente dichiarato, segnalando i disallineamenti.
\item[Tecniche di riferimento] \emph{Using WAI-ARIA role, property, state, and value attributes} (ARIA5/ARIA14/ARIA16); \emph{using standard HTML interactive elements when possible}; \emph{using semantic markup whenever appropriate} (H91).
\item[Condizioni d'uso coperte] Permanenti: ciechi e ipovedenti che usano uno \emph{screen reader} dato che è il criterio che rende possibile l'interazione stessa con qualunque controllo non nativo. Situazionali: navigazione vocale, tecnologie assistive di commutazione (\emph{switch access}).
\end{description}

\paragraph{Criteri presenti in due report (2/3)}

\paragraph{1.3.1 Info and Relationships (Livello A, Percepibile)}
\begin{description}
\item[Verificabile via HTML statico?] Parziale.
\item[Classificazione] Semi-automatica. Molti aspetti sono rilevabili dal solo markup (etichette associate, \texttt{th}/\texttt{scope} nelle tabelle, gerarchia di intestazioni, \texttt{ul}/\texttt{ol} per le liste); capire se il markup rispecchia davvero ciò che è percepibile visivamente, ad esempio testo formattato visivamente come lista ma senza \texttt{ul} o tabelle di layout usate al posto di \texttt{table}, richiede invece di confrontare il rendering con il DOM.
\item[Integrabile con LLM?] Sì: un LLM multimodale può confrontare uno screenshot del rendering con il DOM/markup per individuare discrepanze, sia contenuto che appare strutturato (lista, tabella, intestazione) ma non lo è nel markup, sia markup semantico usato solo per ottenere uno stile visivo (ad esempio una \texttt{table} usata per il layout).
\item[Tecniche di riferimento] \emph{Using h1-h6 to identify headings} (H42); \emph{using \texttt{ol}, \texttt{ul}, \texttt{dl} for lists} (H48); \emph{using table markup to present tabular information} (H51); \emph{using HTML structural elements} (H97/G115); ruoli ARIA per tabelle e landmark (ARIA11-26).
\item[Condizioni d'uso coperte] Permanenti: ciechi e ipovedenti che usano uno \emph{screen reader} e che si affidano alla struttura per navigare; persone con disabilità cognitive che usano la struttura per orientarsi. Situazionali: navigazione vocale; \emph{reflow} su schermi piccoli, che dipende dalla struttura semantica per riorganizzare il contenuto.
\end{description}

\paragraph{1.4.1 Use of Color (Livello A, Percepibile)}
\begin{description}
\item[Verificabile via HTML statico?] Parziale.
\item[Classificazione] Semi-automatica. Alcuni sotto-casi sono calcolabili da CSS/DOM senza screenshot, ad esempio il contrasto fra il colore di un link e il testo circostante, o l'assenza di uno stile o di un'icona aggiuntiva su campi obbligatori segnalati solo tramite \texttt{border-color}. Riconoscere in generale se il colore è l'unico mezzo usato per trasmettere un'informazione (ad esempio la legenda di un grafico, o un'etichetta di stato) richiede invece di vedere il rendering e interpretarne il significato semantico.
\item[Integrabile con LLM?] Sì: un LLM multimodale può analizzare uno screenshot e verificare se un'informazione (stato, categoria, obbligatorietà) è distinguibile anche senza percezione del colore, individuando i casi di campi obbligatori o errori segnalati solo a colore, e di link non distinguibili se non per colore.
\item[Tecniche di riferimento] \emph{Not using color alone to convey information} (G14); \emph{including a text cue whenever color cues are used} (G205); \emph{ensuring additional visual cues are used when color differences are used to convey information} (G182); \emph{using a contrast ratio of 3:1 with surrounding text to distinguish links} (G183); \emph{using color and pattern} (G111).
\item[Condizioni d'uso coperte] Permanenti: persone daltoniche e con altre forme di deficit della percezione cromatica. Situazionali: schermi con calibrazione scarsa o luce solare forte, che riducono la percezione dei contrasti cromatici; dispositivi o modalità in scala di grigi.
\end{description}

\paragraph{2.4.7 Focus Visible (Livello AA, Utilizzabile)}
\begin{description}
\item[Verificabile via HTML statico?] No.
\item[Classificazione] Semi-automatica. Serve un'interazione reale (attivare il focus su ogni elemento) e l'osservazione del rendering per capire se compare un indicatore visibile e distinguibile. Uno strumento può rilevare se esistono regole CSS \texttt{:focus}/\texttt{:focus-visible} definite, o se sono state rimosse senza un sostituto (ad esempio \texttt{outline: none}), ma la sufficienza visiva e di contrasto dell'indicatore richiede l'ispezione del rendering effettivo.
\item[Integrabile con LLM?] Sì, se abbinato ad automazione del browser che catturi lo screenshot con l'elemento a fuoco: un LLM multimodale può confrontare lo stato normale e quello a fuoco di un elemento e giudicare se l'indicatore è visibile e sufficientemente contrastato. Senza automazione che generi gli screenshot per ogni stato di focus, il controllo non è invece applicabile.
\item[Tecniche di riferimento] Tecnica sul contrasto dell'indicatore di focus (C40); \emph{providing a visible focus indicator} (G195); evitare la rimozione dell'\texttt{outline} senza un'alternativa (F78, come \emph{failure} comune).
\item[Condizioni d'uso coperte] Permanenti: disabilità motorie, in particolare chi naviga solo da tastiera; ipovedenti. Situazionali: uso senza mouse, ad esempio con un dispositivo dotato di sola tastiera o di uno \emph{switch}.
\end{description}

\paragraph{3.3.2 Labels or Instructions (Livello A, Comprensibile)}
\begin{description}
\item[Verificabile via HTML statico?] Parziale.
\item[Classificazione] Semi-automatica. La presenza di etichette associate ai campi (tramite \texttt{for}/\texttt{id}, \texttt{aria-label} o \texttt{aria-labelledby}) è rilevabile dal markup; se l'etichetta sia effettivamente comprensibile, completa e adeguata al contesto (ad esempio, istruzioni sul formato atteso per campi complessi) richiede invece un giudizio.
\item[Integrabile con LLM?] Sì: un LLM può giudicare se l'etichetta o l'istruzione associata a un campo è chiara, specifica e sufficiente rispetto a quanto il campo richiede (ad esempio, \enquote{Data} contro \enquote{Data di nascita (gg/mm/aaaa)}), individuando etichette generiche o ambigue e segnalando campi con formati complessi privi di istruzioni.
\item[Tecniche di riferimento] \emph{Providing labels or instructions for user input} (G131); \emph{using \texttt{aria-describedby} to provide a descriptive label} (ARIA1); \emph{identifying required fields and describing them to the user} (G89); \emph{using \texttt{label} elements to associate text labels with form controls} (H44).
\item[Condizioni d'uso coperte] Permanenti: disabilità cognitive; utenti che usano \emph{screen reader}, per cui l'etichetta è l'unica informazione disponibile senza il contesto visivo del modulo. Situazionali: compilazione di moduli in mobilità o in condizioni di distrazione, dove istruzioni chiare riducono gli errori.
\end{description}

\paragraph{4.1.1 Parsing (Livello A, Robusto)}
\textit{Nota:} il criterio 4.1.1 è stato rimosso in WCAG 2.2~\footcite{w3c:wcag22-changes}, che lo ha giudicato ridondante rispetto alla gestione degli errori di parsing ormai garantita dai browser e dalle tecnologie assistive moderne. Questo lavoro fa riferimento a WCAG 2.1 Livello AA (v.~sezione~\ref{sec:livello-aa}), motivo per cui il criterio è comunque incluso in questa analisi preliminare; non sarà tuttavia oggetto di ulteriore sviluppo nei capitoli successivi, essendo già escluso dall'insieme \emph{core} per le ragioni discusse più avanti in questa sezione.
\begin{description}
\item[Verificabile via HTML statico?] Sì.
\item[Classificazione] Automatica. Il criterio richiede che il markup sia ben formato: elementi con tag di apertura e chiusura completi, correttamente annidati, senza attributi duplicati e con \texttt{id} univoci; tutti requisiti verificabili in modo interamente deterministico da un validatore sintattico standard (ad esempio il W3C Markup Validator), senza margini di ambiguità.
\item[Integrabile con LLM?] No, il suo eventuale intervento non aggiungerebbe alcun valore rispetto a un parser tradizionale.
\item[Tecniche di riferimento] \emph{Ensuring that Web pages are well-formed} (G134); uso di un validatore sintattico HTML conforme alle specifiche.
\item[Condizioni d'uso coperte] Permanenti: utenti che usano \emph{screen reader} e di altre tecnologie assistive, il cui corretto funzionamento presuppone un markup analizzabile senza ambiguità dal motore di parsing sottostante.
\end{description}

\paragraph{Criteri presenti in un solo report (1/3)}

\paragraph{1.4.10 Reflow (Livello AA, Percepibile)}
\begin{description}
\item[Verificabile via HTML statico?] Parziale.
\item[Classificazione] Semi-automatica. Alcuni indizi di possibili problemi sono rilevabili attraverso un'analisi statica del CSS, come la presenza di larghezze fisse espresse in unità assolute (\texttt{px}) o l'assenza di \emph{*media query*} per la gestione del layout responsive. Tuttavia, accertare se il contenuto possa essere utilizzato senza dover scorrere la pagina sia orizzontalmente che verticalmente, quando la finestra di visualizzazione ha una larghezza equivalente a 320~CSS px (o al 400\% di zoom), richiede il rendering effettivo della pagina a quella dimensione.
\item[Integrabile con LLM?] Sì, se abbinato ad automazione del browser che renderizzi la pagina alla larghezza equivalente e ne catturi uno screenshot, un LLM multimodale può individuare contenuto tagliato o uno scorrimento orizzontale non necessario. Senza automazione del rendering, il controllo resta limitato ai soli indizi statici nel CSS.
\item[Tecniche di riferimento] \emph{Reflow} (C32, C31, C33, C38); \emph{using percentage values or relative units for container sizing} (G206).
\item[Condizioni d'uso coperte] Permanenti: ipovedenti che utilizzano l'ingrandimento del browser fino al 400\% per leggere il testo. Situazionali: consultazione da dispositivo mobile con schermo di piccole dimensioni, o da una finestra desktop ridotta/affiancata ad altre.
\end{description}

\paragraph{1.4.11 Non-text Contrast (Livello AA, Percepibile)}
\begin{description}
\item[Verificabile via HTML statico?] Parziale.
\item[Classificazione] Semi-automatica, con lo stesso schema già visto per 1.4.3: se i colori di sfondo e degli elementi grafici sono valori CSS solidi dichiarati, il contrasto dei bordi dei componenti dell'interfaccia (campi di input, pulsanti, indicatori di stato) e delle parti grafiche essenziali alla comprensione è calcolabile staticamente. Prima del calcolo è però necessario individuare quali elementi della pagina rientrino effettivamente nell'ambito del criterio, distinguendo i componenti dell'interfaccia e le parti grafiche essenziali alla comprensione dagli elementi che ne sono esclusi; questa valutazione non può essere determinata in modo affidabile dal solo codice CSS.
\item[Integrabile con LLM?] Sì: analogamente a 1.4.3, un LLM multimodale può sia identificare quali elementi dell'interfaccia rientrano nell'ambito del criterio, sia stimare il contrasto su sfondi complessi non calcolabili staticamente, riutilizzando in larga parte la stessa pipeline di analisi già prevista per 1.4.3.
\item[Tecniche di riferimento] G195 (indicatore di focus visibile, esteso al contrasto); G207 (contrasto sufficiente dei bordi dei componenti dell'interfaccia); tecnica analoga a G183, applicata agli elementi non testuali.
\item[Condizioni d'uso coperte] Permanenti: ipovedenti e persone con bassa sensibilità al contrasto, per le stesse patologie oculari già indicate per 1.4.3, qui applicate a icone, controlli di modulo e altri componenti grafici dell'interfaccia. Situazionali: luce solare diretta sullo schermo, schermi di scarsa qualità.
\end{description}

\paragraph{2.4.1 Bypass Blocks (Livello A, Utilizzabile)}
\begin{description}
\item[Verificabile via HTML statico?] Parziale.
\item[Classificazione] Semi-automatica. La presenza di un meccanismo per saltare i blocchi di contenuto ripetuti un collegamento di salto (\emph{skip link}) o un'adeguata struttura di landmark ARIA/HTML5 (\texttt{main}, \texttt{nav}) è rilevabile dal markup; giudicare se il meccanismo è effettivamente sufficiente (landmark che coprono davvero tutte le sezioni ripetute, etichette distintive quando sono presenti più landmark dello stesso tipo) richiede invece un giudizio più fine.
\item[Integrabile con LLM?] Sì: un LLM può analizzare la struttura di landmark e di intestazioni di una pagina e giudicare se copre in modo adeguato e distinguibile i blocchi di contenuto ripetuti, individuando casi limite ad esempio più elementi \texttt{nav} privi di \texttt{aria-label} che li distingua che una semplice verifica di presenza/assenza non coglierebbe.
\item[Tecniche di riferimento] G1 (\emph{skip to main content}), G123, G124; \emph{using landmark roles to identify regions of a page} (ARIA11); \emph{providing heading elements at the beginning of sections} (H69).
\item[Condizioni d'uso coperte] Permanenti: utenti che usano \emph{screen reader} e utenti che navigano esclusivamente da tastiera, per cui saltare i blocchi ripetuti (intestazione, menu di navigazione) ad ogni pagina evita un onere significativo. Situazionali: navigazione vocale.
\end{description}

\paragraph{3.1.1 Language of Page (Livello A, Comprensibile)}
\begin{description}
\item[Verificabile via HTML statico?] Sì.
\item[Classificazione] Automatica. La presenza dell'attributo \texttt{lang} e la validità del relativo tag lingua (ad esempio \texttt{it}, \texttt{en-US}, secondo BCP~47) sono verifiche di solo markup. Resta un sotto-caso semi-automatico: verificare che la lingua dichiarata corrisponda davvero al contenuto reale della pagina.
\item[Integrabile con LLM?] Parziale: il controllo di base non necessita di LLM, trattandosi di puro parsing sintattico. Un LLM è utile solo per il sotto-caso della corrispondenza fra lingua dichiarata e contenuto reale (rilevazione della lingua del testo e confronto con l'attributo \texttt{lang}), un compito semplice ma marginale rispetto al criterio nel suo complesso.
\item[Tecniche di riferimento] \emph{Specifying the human language of Web pages (H57)}.
\item[Condizioni d'uso coperte] Permanenti: utenti che usano \emph{screen reader}, per cui l'attributo serve alla pronuncia corretta da parte del sintetizzatore vocale. Situazionali: utenti udi strumenti di traduzione automatica della pagina.
\end{description}

{\small
\begin{longtable}{|p{1.2cm}|p{3.1cm}|p{1.6cm}|p{1.1cm}|p{1.6cm}|p{1.5cm}|}
\caption{Sintesi dell'analisi dei 13 criteri riscontrati nei report} \label{tab:sintesi-analisi} \\ \multicolumn{6}{p{\linewidth}}{\footnotesize \textit{Nota.} Classificazione: A = automatica, S = semi-automatica, M = manuale. Copertura utenza: P = permanenti, T = temporanee, S = situazionali (solo le condizioni effettivamente coperte). } \\
\hline
\textbf{Crit.} & \textbf{Nome} & \textbf{Verific.} & \textbf{Class.} & \textbf{Integr. LLM} & \textbf{Copert.} \\
\hline
\endfirsthead
\multicolumn{6}{c}{\textit{Tabella~\ref{tab:sintesi-analisi} -- segue dalla pagina precedente}} \\
\hline
\textbf{Crit.} & \textbf{Nome} & \textbf{Verific.} & \textbf{Class.} & \textbf{Integr. LLM} & \textbf{Copert.} \\
\hline
\endhead
\hline
\multicolumn{6}{r}{\textit{continua nella pagina successiva}} \\
\endfoot
\hline
\endlastfoot
\multicolumn{6}{|c|}{\textbf{Presenti in tutti e tre i report (3/3)}} \\
\hline
1.1.1 & Non-text Content & Parziale & S & Sì & P, T, S \\
\hline
1.4.3 & Contrast (Minimum) & Parziale & S & Sì & P, T, S \\
\hline
2.4.4 & Link Purpose (In Context) & Parziale & S & Sì & P, S \\
\hline
4.1.2 & Name, Role, Value & Parziale & S & Sì & P, S \\
\hline
\multicolumn{6}{|l|}{\textbf{Presenti in due report (2/3)}} \\
\hline
1.3.1 & Info and Relationships & Parziale & S & Sì & P, S \\
\hline
1.4.1 & Use of Color & Parziale & S & Sì & P, S \\
\hline
2.4.7 & Focus Visible & No & S & Sì & P, S \\
\hline
3.3.2 & Labels or Instructions & Parziale & S & Sì & P, S \\
\hline
4.1.1 & Parsing & Sì & A & No & P \\
\hline
\multicolumn{6}{|l|}{\textbf{Presenti in un solo report (1/3)}} \\
\hline
1.4.10 & Reflow & Parziale & S & Sì & P, S \\
\hline
1.4.11 & Non-text Contrast & Parziale & S & Sì & P, S \\
\hline
2.4.1 & Bypass Blocks & Parziale & S & Sì & P, S \\
\hline
3.1.1 & Language of Page & Sì & A & Parziale & P, S \\
\hline
\end{longtable}
}

\subsection*{Selezione finale: criteri core, opzionali ed esclusi}
 
Incrociando i risultati dell'analisi con i criteri di selezione introdotti in apertura di sezione -- frequenza, integrabilità con LLM e copertura delle condizioni d'uso -- a cui si aggiunge, per i criteri con dati disponibili in entrambe le rilevazioni AGID, l'ampiezza del divario fra la quota di errore rilevata automaticamente e quella confermata dal controllo manuale discussa sopra, i 13 criteri riscontrati nei report sono stati infine suddivisi in tre gruppi: \emph{core}, \emph{opzionali} ed \emph{esclusi}.
 
Sono stati selezionati come \emph{core} otto criteri, riassunti nella tabella~\ref{tab:criteri-core}. Rispetto ai criteri di selezione già introdotti nella sezione~\ref{sec:livello-aa}, due considerazioni aggiuntive hanno avuto un ruolo nella scelta finale. In primo luogo, il confronto fra le percentuali del monitoraggio semplificato e approfondito ha motivato l'inclusione di 1.3.1 (Info and Relationships) fra i criteri \emph{core}: come illustrato sopra, è il criterio con il divario più ampio fra le due rilevazioni (dall'1,6\% al 23,2\%) e con la quota più alta in assoluto nel monitoraggio approfondito -- un'evidenza diretta del fatto che si tratta di un errore sistematicamente sotto-rilevato dal controllo automatico e sistematicamente individuato dal controllo manuale, esattamente il profilo di criterio per cui un sistema assistito da LLM promette il maggiore valore aggiunto. Resta comunque valida l'osservazione che 1.3.1 richiede di essere scomposto in numerosi sotto-casi (intestazioni, liste, tabelle, relazioni implicite fra elementi visivamente correlati): il maggiore costo di implementazione e validazione che ne deriva, discusso nei capitoli~\ref{chap:progettazione-prototipo} e~\ref{chap:validazione-benchmark}, è considerato giustificato dalla rilevanza empirica del criterio.
 
In secondo luogo, lo stesso principio è stato applicato a 2.4.1 (Bypass Blocks), l'unico criterio, fra quelli riscontrati, presente esclusivamente nel monitoraggio approfondito e non nei due report basati su motori automatici (v.~tabella~\ref{tab:frequenza-errori}): il fatto che emerga soltanto dalla valutazione manuale è di per sé un indizio che si tratti di un aspetto poco accessibile all'analisi puramente sintattica, e l'analisi condotta sopra conferma che il criterio è effettivamente integrabile con un LLM, che può giudicare l'adeguatezza della struttura di landmark rispetto ai blocchi di contenuto ripetuti. Per i restanti criteri con entrambi i valori AGID disponibili, il confronto fra le due percentuali non ha invece motivato ulteriori inclusioni: l'ampiezza del divario risulta già più marcata proprio nei tre criteri appena discussi (1.3.1, ma anche 4.1.2 e 2.4.4, già inclusi per frequenza), mentre per i criteri con un divario di segno opposto, come 2.4.7, tale confronto non è stato considerato un argomento a favore dell'inclusione, per le ragioni già discusse sopra. Va inoltre segnalato che il confronto fra le due percentuali non è sempre applicabile in modo pienamente rigoroso, derivando da campioni e metodologie di rilevazione differenti (motore automatico su larga scala contro valutazione manuale su un campione più ristretto): per i criteri per cui uno dei due valori non è disponibile, come lo stesso 2.4.1 o 3.3.2, il confronto diretto non è stato quindi effettuato, e la motivazione di inclusione o esclusione si è basata sugli altri criteri di selezione.
 
Il criterio 2.1.1 (Keyboard), pur restando l'unico a coprire in modo specifico le disabilità motorie che impediscono l'uso del mouse, non compare fra le prime 10 posizioni di nessuno dei due cruscotti AGID aggiornati, e non è quindi presente in alcuno dei tre report di riferimento assunti come base per la selezione: coerentemente con la scelta di basare la selezione unicamente sui dati aggiornati, il criterio non è stato incluso fra i criteri riscontrati né, di conseguenza, fra i criteri \emph{core}. Lo stesso vale per 3.2.2 (On Input) e 4.1.3 (Status Messages), anch'essi assenti dalle prime 10 posizioni dei cruscotti aggiornati.
 
{\small
\begin{longtable}{|p{1.4cm}|p{2.6cm}|p{2.1cm}|p{1.8cm}|p{3.9cm}|}
\caption{Criteri di successo selezionati per la verifica sperimentale (insieme \emph{core}).}
\label{tab:criteri-core} \\
\hline
\textbf{Criterio} & \textbf{Nome} & \textbf{Principio} & \textbf{Classific.} & \textbf{Motivazione sintetica} \\
\hline
\endfirsthead
\multicolumn{5}{l}{\textit{Tabella~\ref{tab:criteri-core} -- segue dalla pagina precedente}} \\
\hline
\textbf{Criterio} & \textbf{Nome} & \textbf{Principio} & \textbf{Classific.} & \textbf{Motivazione sintetica} \\
\hline
\endhead
\hline
\multicolumn{5}{r}{\textit{continua nella pagina successiva}} \\
\endfoot
\hline
\endlastfoot
1.1.1 & Non-text Content & Percepibile & Semi-automatica & Frequenza altissima; copertura utenza completa; caso più ricco di integrazione LLM (confronto visivo-testuale e inferenza del messaggio funzionale). \\
\hline
1.4.3 & Contrast (Minimum) & Percepibile & Semi-automatica & Frequenza altissima; quasi interamente calcolabile con formula deterministica, funge da controllo metodologico rispetto ai criteri valutati dall'LLM. \\
\hline
1.3.1 & Info and Relationships & Percepibile & Semi-automatica & Divario più marcato fra tutti i criteri dal controllo automatico (1,6\%) a quello manuale (23,2\%, il valore più alto del monitoraggio approfondito); massimo valore aggiunto atteso da un giudizio semantico assistito da LLM, nonostante il maggior costo di scomposizione in sotto-casi. \\
\hline
4.1.2 & Name, Role, Value & Robusto & Semi-automatica & Frequenza alta, in forte crescita dall'automatico (1,1\%) al manuale (17,8\%); criterio cardine per l'interazione via screen reader con qualunque widget non nativo. \\
\hline
2.4.4 & Link Purpose (In Context) & Utilizzabile & Semi-automatica & Frequenza alta, in crescita dall'automatico (2,3\%) al manuale (11,1\%); compito puramente testuale per l'LLM, basso costo implementativo. \\
\hline
3.3.2 & Labels or Instructions & Comprensibile & Semi-automatica & Frequenza alta; stesso pattern LLM di 1.1.1 applicato ai moduli, ne dimostra la generalizzabilità. \\
\hline
2.4.7 & Focus Visible & Utilizzabile & Semi-automatica & Criterio con la quota di errore più alta in assoluto fra tutti quelli rilevati (42,4\% nel monitoraggio semplificato AGID); il forte divario con il manuale (2,6\%) resta un'ipotesi da verificare in sede di validazione, non ancora una causa accertata. \\
\hline
2.4.1 & Bypass Blocks & Utilizzabile & Semi-automatica & Unico criterio presente esclusivamente nel monitoraggio approfondito (5,1\%), assente dai report basati su motori automatici; integrabile con LLM, colma un'ulteriore lacuna dell'automazione. \\
\hline
\end{longtable}
}
 
L'insieme selezionato copre tutti e quattro i principi POUR (v.~sezione~\ref{sec:quattro-principi}): Percepibile (1.1.1, 1.4.3, 1.3.1), Utilizzabile (2.4.1, 2.4.4, 2.4.7), Comprensibile (3.3.2) e Robusto (4.1.2), con una rappresentazione del principio Percepibile ora più solida rispetto a quanto indicato nella sezione~\ref{sec:livello-aa}; il principio Utilizzabile, privo del contro-esempio motorio rappresentato da 2.1.1, resta comunque coperto da tre criteri.
 
Un secondo gruppo di criteri è stato classificato come \emph{opzionale}: 1.4.1 (Use of Color), 1.4.10 (Reflow) e 1.4.11 (Non-text Contrast). Si tratta, in ciascun caso, di criteri appartenenti alla stessa famiglia tecnica di uno dei criteri \emph{core} già selezionati (rispettivamente: percezione visiva del colore, rendering responsivo affine a 2.4.7, calcolo del contrasto affine a 1.4.3), ma che non introducono un pattern di verifica sostanzialmente nuovo rispetto ad esso.
 
Sono stati infine \emph{esclusi} due criteri, in entrambi i casi perché la loro verifica è quasi interamente riconducibile a un controllo sintattico tradizionale, per cui l'apporto di un LLM risulterebbe marginale: 3.1.1 (Language of Page), la cui validità si riduce alla presenza e correttezza formale dell'attributo \texttt{lang}, e 4.1.1 (Parsing), interamente verificabile da un validatore sintattico HTML standard. Per quest'ultimo, l'esclusione è ulteriormente motivata dal fatto che il criterio è stato rimosso in WCAG 2.2~\footcite{w3c:wcag22-changes}: pur essendo parte del riferimento normativo WCAG 2.1 Livello AA adottato in questo lavoro, non verrà quindi analizzato oltre questa sezione. Entrambi i criteri sono comunque richiamati in questo lavoro come termine di paragone, a evidenziare che non tutti i criteri traggono lo stesso beneficio da un'architettura basata su LLM.
 
La classificazione completa dei 50 criteri, comprensiva delle motivazioni di dettaglio e delle tecniche di riferimento W3C associate a ciascuno, è riportata integralmente nel materiale di lavoro dello stage e viene ripresa, per i soli criteri selezionati come \emph{core}, nell'analisi dei requisiti del capitolo~\ref{chap:progettazione-prototipo}.
 
Il numero di otto criteri \emph{core} non deriva unicamente dai criteri di selezione discussi sopra, ma tiene conto anche delle tempistiche effettivamente disponibili per la loro implementazione, previste dal piano di lavoro dello stage. Le settimane 2--4 coprono, rispettivamente, la definizione della strategia di prompting e la realizzazione del primo prototipo con un LLM di riferimento (\emph{frontier model}), il lavoro di prompt engineering e la sperimentazione con modelli di dimensioni minori a fini di confronto, e infine la validazione sperimentale, comprensiva dell'analisi di falsi positivi, falsi negativi e allucinazioni: un arco di tre settimane entro cui va quindi realizzata un'estensione web funzionante con integrazione LLM, testata su più tentativi e più modelli, e infine validata su un benchmark di casi reali. Un insieme di criteri troppo ampio avrebbe compromesso la possibilità di condurre, per ciascun criterio, un numero di iterazioni di prompt engineering e un confronto fra modelli sufficienti a produrre risultati significativi in questo intervallo di tempo, oltre a ridurre il tempo disponibile per la validazione stessa.
 
Nella scelta ha pesato anche la complessità implementativa dei singoli criteri, non solo il loro numero. Alcuni dei criteri \emph{core}, come 1.3.1, richiedono la gestione di numerosi sotto-casi (intestazioni, liste, tabelle, relazioni implicite fra elementi) e quindi un prompt engineering più articolato; altri, come 2.4.7, presuppongono l'acquisizione automatizzata dello stato visivo della pagina in condizioni diverse (ad esempio schermate a fuoco spostato fra i vari elementi interattivi), il che richiede funzionalità di automazione del browser aggiuntive rispetto alla sola interrogazione dell'LLM sul markup statico; altri ancora, come 2.4.4 o 3.3.2, si prestano invece a un controllo prevalentemente testuale, con un costo di implementazione contenuto. L'insieme finale di otto criteri è stato quindi calibrato per restare realizzabile entro le tre settimane previste, includendo comunque una rappresentanza di criteri a complessità implementativa eterogenea, in modo da non limitare la validazione ai soli casi più semplici.
 
È importante distinguere questa valutazione di fattibilità dai criteri di selezione discussi in precedenza: la rilevanza empirica di un criterio (frequenza, divario fra le rilevazioni, copertura d'utenza) resta il fattore che ne determina l'inclusione nell'insieme \emph{core}, indipendentemente dal costo implementativo. In particolare, 2.4.7 e 1.1.1 -- inclusi per la frequenza più alta in assoluto e per frequenza e copertura d'utenza massime, rispettivamente -- presentano anche il rischio implementativo più elevato dell'insieme: 2.4.7 richiede di dare il focus programmaticamente a ciascun elemento interattivo della pagina e di confrontarne lo stato visivo risultante, funzionalità di automazione del browser aggiuntive rispetto alla sola interrogazione dell'LLM sul markup statico; 1.1.1 richiede, per un giudizio accurato sull'adeguatezza di un testo alternativo, la capacità di mettere in relazione il contenuto visivo di un'immagine con il testo che la descrive, funzione che presuppone un LLM multimodale o un passaggio dedicato di descrizione dell'immagine. Anziché escludere questi due criteri dal core in base al costo di implementazione, si è scelto di mantenerne la classificazione motivata dalla rilevanza empirica e di prevedere, nella progettazione del prototipo (capitolo~\ref{chap:progettazione-prototipo}), una versione semplificata per ciascuno, da estendere compatibilmente con i tempi a disposizione: per 2.4.7, un controllo statico preliminare basato sui fogli di stile (ad esempio la rilevazione di regole che rimuovono esplicitamente l'indicatore di focus di default, come \texttt{outline: none} senza un sostituto visibile), a cui affiancare l'automazione per-elemento se il tempo lo consente; per 1.1.1, un'euristica basata sul solo testo (confronto fra nome del file, attributi presenti nel markup e testo alternativo dichiarato), a cui affiancare la valutazione multimodale come estensione.
 
\FloatBarrier